\section{Proof of Theorem~\ref{th:edit}}\label{app:proof}
We count changes to the \(D=d^2+d\) real entries of the transition matrix \(A\) and input column \(b\) of a width-\(d\) linear recurrent network. The output row \(c\) reads one fixed state coordinate. Renumbering the coordinates preserves edit counts, so take \(c=e_1^\top\), where \(e_1\) is the first standard basis column and \({}^\top\) denotes transpose. Here, the real numbers form the field of coefficients for every calculation in the proof. A unit impulse is input 1 on the first tick and 0 thereafter. Index its desired scalar outputs by \(g_0,g_1,\ldots\) and fix the number \(H\ge1\) of required outputs. For a positive integer \(s\), define the \emph{observation matrix} \(O_s(A)\) by stacking the \(s\) rows \(c,cA,\ldots,cA^{s-1}\). The rank of a matrix is the number of independent columns (equivalently, rows). Its kernel comprises the vectors that it sends to zero. The product of \(O_s(A)\) with \(b\) gives the network's first \(s\) impulse outputs. Write \(g_{0:s-1}=(g_0,\ldots,g_{s-1})^\top\).

Let \(S\) be the forward shift on sequences, so that \((Sg)_i=g_{i+1}\). A monic polynomial has leading coefficient one. For a monic polynomial of degree \(m\), written \(q(z)=z^m+\sum_{j=0}^{m-1}q_jz^j\) in the indeterminate \(z\), set
\[
             (q(S)g)_i=g_{i+m}+\sum_{j=0}^{m-1}q_jg_{i+j}.
\]
Fix the smallest-degree monic polynomial \(p\) for which \(p(S)g=0\) at all indices. Its degree \(r\) is the recurrence order of the nonzero target. We assume \(d\ge r\) until the final step. A \emph{Hankel matrix} is constant on each antidiagonal. The leading \(r\)-by-\(r\) Hankel matrix \(G_r=(g_{i+j})_{0\le i,j<r}\) is invertible. Otherwise a nonzero linear combination of its columns would give a nonzero polynomial of degree less than \(r\) whose filtered sequence has \(r\) initial zeros. Filtering commutes with \(p(S)\), so the recurrence propagates those zeros to every index, contradicting minimality. All columns of the infinite Hankel matrix after the first \(r\) obey the same recurrence, giving overall rank \(r\).

\paragraph{Controlling the characteristic polynomial.}
Call \((c,A)\) \emph{observable} if \(O_d(A)\) has rank \(d\), so the first \(d\) scalar observations distinguish all initial states. Let \(I\) be the \(d\)-by-\(d\) identity, let \(\chi_A(z)=\det(zI-A)\) be the characteristic polynomial, and let \(\mathrm{adj}(M)\) denote the transpose of the cofactor matrix of a square matrix \(M\). For a freely chosen column \(u\in\mathbb R^d\), changing \(A\) to \(A+uc\) changes only column one. The rank-one determinant identity gives
\[
         \chi_{A+uc}(z)-\chi_A(z)=-c\,\mathrm{adj}(zI-A)u .
\]
For observable \((c,A)\), the displayed map from \(u\) to polynomials of degree below \(d\) is a linear bijection. Indeed, if the right side vanishes, so does \(c(zI-A)^{-1}u\) wherever \(zI-A\) is invertible. Its power series in \(1/z\) near infinity has coefficients \(cA^ku\) for integers \(k\ge0\). They all vanish, hence \(O_d(A)u=0\) and \(u=0\). Both spaces have dimension \(d\), proving bijectivity. The new matrix remains observable because \(c(A+uc)^k\) equals \(cA^k\) plus a linear combination of preceding observation rows. Finally, the determinant of \(O_d(A)\) is a polynomial which equals one for the matrix with ones on its first superdiagonal and zeros elsewhere. Its zero set has Lebesgue measure zero. Thus the required observability holds with probability one under the initialization in Theorem~\ref{th:edit}.

\paragraph{Constructing an edit.}
If \(H\le d\), select \(H\) linearly independent columns of \(O_H(A)\). Holding the remaining entries of \(b\) fixed, solve the resulting invertible system for the selected entries of the new input vector \(b'\) to obtain \(O_H(A)b'=g_{0:H-1}\). This takes at most \(H\) edits.

For \(H>d\), put \(k=H-d\). Seek a monic degree-\(d\) polynomial \(q\) such that the first \(k\) terms of \(q(S)g\) vanish. These linear equations in the coefficients of \(q\) are consistent because a monic degree-\(d\) multiple of \(p\) exists. Their \(k\)-by-\(d\) Hankel coefficient matrix has rank \(\min(k,r)\), since it contains the first \(\min(k,r)\) rows of the invertible \(G_r\) and the infinite matrix has rank \(r\). The bijection above lets us impose the equations on \(q=\chi_{A+uc}\) with the same rank in \(u\). Select that many independent columns of this linear system and use just their coordinates of \(u\). Then solve \(O_d(A+uc)b'=g_{0:d-1}\), using at most \(d\) more edits. The Cayley--Hamilton identity says a square matrix satisfies its characteristic polynomial. Consequently the network impulse outputs satisfy the recurrence \(q\), and so agree by induction with the target through the first \(H\) outputs. If \(k\ge r\), the first \(r\) terms of \(q(S)g\) vanish, and the recurrence \(p\) propagates those zeros to all subsequent terms. The total cost is at most \(d+\min(k,r)\).

\paragraph{Bounding the size of the solution set.}
Codimension is the ambient dimension minus the dimension of the subset. We show that the set of \((A,b)\) realizing the prefix has codimension at least \(\min(H,d+r)\), allowing for unobservable realizations. Let \(m=\rank O_d(A)\), and write \(U=\ker O_d(A)\) for the subspace of initial states invisible to the output. Because \(c\ne0\), we have \(1\le m\le d\). The subspace \(U\) has dimension \(d-m\), lies in \(\ker c\), and satisfies \(AU\subseteq U\) by Cayley--Hamilton. The family of possible subspaces \(U\) of this dimension inside the \((d-1)\)-dimensional space \(\ker c\) has dimension \((d-m)(m-1)\). Locally such a subspace is the graph of a linear map between spaces of dimensions \(d-m\) and \(m-1\). For fixed \(U\), the condition \(AU\subseteq U\) forces a block of \(m(d-m)\) entries to vanish in a basis adapted to \(U\). Subtracting the freedom to choose \(U\) leaves codimension at least \(d-m\). Identifying states that differ by an element of \(U\) gives an \(m\)-dimensional quotient space. Denote its induced transition and output by \(\bar A,\bar c\). This quotient is observable.

If \(m\ge H\), the first \(H\) observation rows are independent. Once an observation row belongs to the span of earlier rows, multiplying the relation by \(A\) prevents every subsequent row from increasing that span. Hence the \(H\) impulse equations impose \(H\) constraints on \(b\) for fixed \(A\), as required. If \(m<H\), the first \(m\) impulse equations impose \(m\) independent constraints on \(b\), leaving \(d-m\) free coordinates. Put \(k=H-m\). The other equations require \(q=\chi_{\bar A}\) to satisfy \((q(S)g)_i=0\) for \(0\le i<k\). The coefficient matrix \(B=(g_{i+j})_{0\le i<k,\,0\le j<m}\) has rank \(\rank B\). When consistent, these equations impose that many independent conditions on \(\bar A\). To see this, perturb it by \(\bar u\bar c\) for a quotient-space column \(\bar u\). The characteristic coefficients vary freely by the bijection above, and every such perturbation lifts to a matrix that preserves \(U\). Thus the total codimension for this value of \(m\) is at least \((d-m)+m+\rank B=d+\rank B\).

If \(m\ge r\), then \(\rank B=\min(k,r)\) and \(d+\rank B\ge\min(H,d+r)\). If \(m<r\) and \(k\ge r\), propagation of the initial zeros of \(q(S)g\) would give an annihilating polynomial of degree \(m<r\), contradicting the definition of \(r\). Finally if \(m<r\) and \(k<r\), deleting \(r-k\) rows and \(r-m\) columns of \(G_r\) leaves \(B\), so \(\rank B\ge k+m-r=H-r\). Then \(d+\rank B\ge H\), as required.

\paragraph{From dimension to scalar edits.}
Fix a set of \(s\) coordinates allowed to change. Initial networks which can reach the prefix solution set lie in its sum with the \(s\)-dimensional subspace supported on those coordinates. The solution set is defined by polynomial equations in real parameters. Its sum with this subspace is semialgebraic, meaning that it can be described using finitely many real polynomial equalities and inequalities and Boolean combinations. Its dimension is at most \(D-\min(H,d+r)+s\) because it is the image under addition of a product of that dimension at most. If \(s<\min(H,d+r)\), this is less than \(D\), so the sum has Lebesgue measure zero. There are finitely many choices of coordinates and countably many positive integers \(H\), giving the lower bounds simultaneously with probability one. At \(H=d+r\), the upper and lower constructions give the all-horizon result. For \(d<r\), Cayley--Hamilton would give an annihilating polynomial of degree \(d\), ruling out all-horizon realization. \hfill\(\square\)

\paragraph{Exact arithmetic verification.}
We treated each binary64 Xavier coefficient as its exact rational value. The three test impulse sequences were the constant \(g_n=1\), the Fibonacci numbers \(F_0=0,F_1=1,F_{n+2}=F_{n+1}+F_n\), and the triangular numbers \(T_n=n(n+1)/2\) for nonnegative integers \(n\). We tested five initialization seeds and widths \(1,2,3,4,6,8\) whenever \(d\ge r\). Across 480 cases with \(1\le H\le d+r\), every write changed exactly the predicted number of coordinates. Each of the 405 shorter-horizon writes failed at the next coefficient \(g_H\). At \(H=d+r\), all 75 characteristic polynomials of the edited matrices were divisible by \(p\), and their first 32 impulse coefficients matched directly. Divisibility and agreement on the first \(d\) coefficients then certify every later coefficient.

\section{A comparator that preserves binary64 operands}\label{app:sort}
A comparator takes two inputs and returns their smaller and larger values in that order. Let \(a,b\) be finite binary64 inputs in \([-10,10]\), let \(R(t)=\max(t,0)\) be the rectified linear unit, and take fixed constants \(K=2^{537}\), \(M=10\). Round each displayed operation separately to the nearest binary64 value, choosing the one with even last significand bit at a tie. Gradual underflow retains reduced-precision values below the smallest positive normal number. Numerical equality identifies the two signed zeros.

First extract the positive difference \(r\), amplify it to \(u\), and compute the auxiliary signals \(v,B,S\). Then compute the keep and swap selectors \(k,s\), respectively, in this order:
\[
 r=R(a-b),\quad u=Kr,\quad v=u-R(u-1),\quad B=1-R(1-u),
 \quad S=1-R(1-Kv),
\]
\[
                k=R((1-B)-S),\qquad s=1-k.
\]
For an operand \(x\) and an enable bit \(g\in\{0,1\}\), set
\(G(x,g)=R(x-M(1-g))-R(-x-M(1-g))\). The comparator returns
\((G(a,k)+G(b,s),\,G(a,s)+G(b,k))\), reusing the shifts \(Mk,Ms\).

When \(a\le b\), \(u=B=S=0\) and \(k=1\). When \(a>b\), the rounded positive difference is at least the smallest representable positive value \(2^{-1074}\). For \(0<u\le1\), \(v=u\) and \(Kv\ge1\), hence \(S=1\). For \(u\ge1\), \(B=1\). In that case the rounded value of \(u-1\) lies between \(u-2\) and \(u\), since \(u\) itself is representable and is one unit of absolute distance from the exact difference. Therefore \(0\le v\le2\), which also shows that the product \(Kv\) remains finite. Both \(B,S\) lie in \([0,1]\) and at least one equals one, so \(k=0\). For \(g=1\) the gate \(G\) returns the original operand, while for \(g=0\) both rectifiers vanish. Composition along a sorting schedule therefore preserves all numerical operands. The stored comparator contains 43 edge weights and seven biases, with Xavier initialization for the weights and zero initialization for the biases.

\section{Scalar construction and size sweeps}\label{app:square}
For a real input \(x\in[-16,16]\), set \(t=|x|/16\). Define the tent map \(T(t)=2t-4R(t-1/2)\) on \([0,1]\), where \(R(t)=\max(t,0)\). For a positive integer refinement depth \(m\), let \(T^{\circ i}\) denote \(i\) successive applications of \(T\). The depth-\(m\) square interpolant is
\[
               256\left[t-\sum_{i=1}^m4^{-i}T^{\circ i}(t)\right].
\]
At each refinement, the added term corrects the midpoint error of the preceding linear interpolant on intervals with endpoints at multiples of a power of two. The maximum absolute error of the resulting interpolant on \([-16,16]\) in real arithmetic is \(64\cdot4^{-m}\). The compact model starts at hidden width two, uses width three in its remaining layers, and has an affine output and \(12m+5\) parameter slots.

For the compensated version let \(q\) denote the running unscaled sum and \(u\) its next term to subtract. Compute the rounded difference \(s=q-u\), then \(d=q-s\), then \(e=d-u\), using the compensated subtraction transform \cite{dekker1971}. Accumulate \(R(e)\) and \(R(-e)\) separately. Add their difference to the last \(q\) before scaling the output. Copying nonnegative intermediate results through rectifiers implements the stated arithmetic with maximum hidden width six and \(203m+35\) slots. In exact real arithmetic \(e=0\). Setting only the two weights that read the error accumulators to zero tests the numerical contribution of the compensation.

To construct difficult shell inputs, express a positive normal input as \(x=N\,2^{e-52}\), where the integer exponent \(e\) places \(x\) in the interval \([2^e,2^{e+1})\) and \(N\) is its 53-bit integer significand. The low bits of \(N^2\) determine which side of a rounding midpoint contains \(x^2\). Let \(k\in\{52,53\}\) count the discarded bits and take \(j\in\{7,15,\ldots,255\}\). Find integers satisfying \(N^2\equiv2^{k-1}-j\pmod{2^k}\) by extending the low bits one at a time. Retain significands with the specified discard count and shell magnitude, then include both signs. This gives 304 inputs from 152 magnitudes. Networks with 35, 48 and 55 refinements match respectively \(0,40,304\) of these inputs. All three match all 131,072 fresh uniform shell inputs. On 501 small powers \(x=2^n\), \(-500\le n\le0\), the same networks match 32, 45 and 52 inputs. The first affine segment near zero gives large relative errors despite small absolute errors.
